\section{Discussion}
\label{sec:discussion}

\paragraph*{Hypothesis verdicts at a glance.}
\Cref{tab:hypothesis-verdicts} lets us read the structure of the
findings as a small number of compatible patterns rather than as a long
list of numbers. Cost is heterogeneous (H1a) and selection-biased by
strategy (H1b) wherever $\TrainCost$ is non-trivial, and on those same
scopes $n$ is well rank-correlated (H2a) but strategy-biased (H2c).
H2b --- ``the median calibration ratio is not centered at $1$'' --- is
supported only when the practical bias clears a $5\%$ floor; on the
remaining systems the ratio is statistically off-$1$ but not
practically. The decision-level claims H3a, H4a and H4b are existential
(``rankings sometimes diverge''/``some pairs flip Pareto status'')
and they hold as soon as a single scope flips, which happens on
$3/7$ \StudyReplication systems for H3a, $5/7$ for H4a/H4b, and on
$25$--$65/68$ x264 input/scenario pairs. Cross-system, the
strongest signal is H3b: rank-flip count and per-configuration cost CV
correlate at Spearman $\rho=0.90$ on seven systems, $p=6.3\!\times\!
10^{-3}$. Heterogeneity is therefore the single best predictor of
cost-model fragility, and the diagnostic below operationalizes that
finding.

\paragraph*{When is $n$ safe?}
The empirical answer is system-, scenario-, and context-dependent. On
the \StudyReplication batch, four systems (Dune, HIPAcc, LLVM, Polly) support the
size-as-cost abstraction without flipping any top-$1$ recommendation
across $\lambda$ values; for those systems sample size is a usable cost
proxy in practice. Three (7Z, BerkeleyDBC, and especially lrzip) do not:
their per-configuration cost CV exceeds $1$, calibration ratios drift
$10$--$20\%$ from $1$, and rankings flip in $\geq\!1/7$ budgets. On
x264, even when calibration is benign in aggregate, the
cost-target scenario matters: predicting size while paying time-based
cost ($\text{cost}=\text{time}, \text{target}=\text{size}$) is the
weakest scenario in both rank agreement ($0.82$) and Pareto overlap
($0.38$).

\paragraph*{When does empirical cost matter most?}
Two predictors of decision-level disagreement emerge. The first is
heterogeneity: lrzip ($p_{90}/p_{10}\!=\!47.0$, $\mathrm{CV}\!=\!1.98$)
has the most rank flips ($5/7$ across $\lambda$), and the cross-system
H3b correlation between cost CV and rank-flip count is Spearman
$\rho\!=\!0.90$ ($p\!=\!6.3\!\times\!10^{-3}$, $n\!=\!7$ \StudyReplication
systems). Pareto-front instability is more subtle: the lowest Pareto
Jaccard is on LLVM ($0.33$), not on the most heterogeneous systems,
because under count cost LLVM admits only three non-dominated pairs
while empirical cost reveals nine; this is a low-CV system where the
empirical-cost view simply has more discriminative power than the
count view, not a fragility induced by heterogeneity. The second
predictor is target--cost asymmetry: predicting size while paying
time-based cost on x264 is the worst scenario despite x264's
relatively benign cost CV ($0.21$--$0.39$). A practitioner whose
acquisition cost is dominated by a single resource and whose
prediction target differs from that resource is exactly the case
where empirical cost should not be proxied by $n$.

\paragraph*{A practical diagnostic.}
The cost CV and the calibration ratio
$\mathit{TrainCostPct}/\mathit{SizePct}$ can both be estimated from a
small pilot sample. We recommend the following lightweight check before
committing to a size-based budget: (i) measure a uniform random pilot
of $\sim\!\!2$--$5\%$ of $|C|$ and compute the cost CV, (ii) for each
candidate strategy, draw a few small samples and compute the median
$\mathit{TrainCostPct}/\mathit{SizePct}$, (iii) if either CV $>0.5$ or
the per-strategy ratio is far from $1$, switch the budget objective to
$\TrainCost$ rather than $n$. The cost reported as RQ2 calibration in
the \StudyReplication systems can be reproduced from a $5\%$ pilot in seconds for
all systems we studied.

\paragraph*{Linux complements but does not aggregate.}
The Linux cross-version study (\StudyLinux family in
\Cref{sec:results}) is reported separately because its protocol
differs structurally from \StudyReplication: pool-based subsampling instead of synthesis-based sampling,
storage-bound cost ($\mathit{cost}(c)$ $=$ uncompressed binary size)
instead of time-based cost, and a strategy set that includes
\texttt{uncertainty} (active learning) and \texttt{diversity}
(farthest-first) which have no analogues in the \StudyReplication
protocol. The clean headline on Linux is neither ``always use empirical
cost'' nor ``$n$ is a fine proxy'': across seven releases
(v4.13, v4.15, v4.20, v5.0, v5.04, v5.07, v5.08), \emph{calibration is a property of the
strategy, not the system}. Empirical $\TrainCost$ is required as
soon as the strategy panel includes AL or diversity-aware samplers
(\texttt{diversity} pays $\sim\!\!2\times$ its count budget on
\emph{every} release: medians $2.28, 2.23, 2.27, 2.06, 2.09, 2.00, 2.02$ (window $[2.00, 2.28]$) across seven\nindependent randconfig pools, the most robust positive finding in
the paper; \texttt{uncertainty}'s drift direction is sensitive to
absolute $n$, with cross-over $n$ pool-dependent); $n$ remains a
usable proxy for cost-blind samplers (\texttt{random},
\texttt{balanced}, \texttt{coverage} are calibrated within $\\pm 0.7\\%$ of $1$ on every release with Spearman $\rho\!\geq\!0.97$).
Three of the synthesis-based findings travel: (i) heavy-tailed cost
is again the most reliable predictor of decision-level instability
\emph{within a fixed protocol} (Linux v4.13 CV $1.43$ behaves like
lrzip CV $1.98$), but the H3b cross-system correlation does
\emph{not} extend to the Linux pool: three of seven \StudyLinux releases
share the same $2/7$ flip count across pool sizes $19{,}952$--$92{,}562$ and CVs $0.99$--$1.43$, while v5.04 flips $7/7$, v5.07 flips $4/7$, and v5.08 flips $6/7$ for a
level-grid reason unrelated to its CV; (ii) the empirical-cost
Pareto front is again strictly larger than the count-cost front on
all seven releases; (iii) the calibration story is again
strategy-specific rather than a single per-system aggregate. Three
findings are Linux-specific: \texttt{diversity} drifts to
$\sim\!\!2\times$ expected cost on all seven releases (7$\times$ replication), \texttt{uncertainty} exhibits an $n$-dependent
calibration seam whose direction is sensitive to absolute $n$, and
ranking instability tracks the level grid rather than pool size or
cost CV (v5.0, the smallest pool, flips $2/7$ like the larger v4.13
and v4.15; v5.04, v5.07, v5.08 with their tie-prone level grids flip $7/7$, $4/7$, $6/7$).
None of these has an analogue in the synthesis-based studies
because the corresponding strategies and pool-based regime are not
in the synthesis-based panel; this is a reason to keep Linux
separate, not to aggregate.

\paragraph*{\StudyEnergy complements but does not aggregate.}
The the performance--energy corpus generated-sampling study
(\StudyEnergy family in \Cref{sec:results}) is reported separately for the same
reason as Linux: the protocol differs structurally. \StudyEnergy draws
samples from a measured configuration universe under a strategy panel
(\texttt{random}, \texttt{balanced}, \texttt{diversity},
\texttt{pairwise}, \texttt{pilotCost}) that overlaps neither with the
synthesis-based panel nor with the Linux pool-based panel, and
exposes two cost properties (\emph{performance} and \emph{energy})
which let us probe a non-time cost dimension explicitly. The
qualitative findings travel: $n$ is rank-correlated with empirical
cost ($\rho\!\geq\!0.71$ on every system), strategy-specific
calibration bias survives when aggregate medians look calibrated, and
empirical-cost Pareto frontiers differ from count frontiers in
$28/28$ scenarios. Quantitatively, \StudyEnergy also extends the dynamic
range of the paper: jump3r is the most heterogeneous case anywhere in
this work (CV $\!=\!7.10$), and PostgreSQL\,11.2 is the most
homogeneous (CV $\!=\!0.06$); both exist within the same protocol,
which is informative for any practitioner whose acquisition cost is
in the ``unknown CV'' regime. the \StudyEnergy specific findings include the
stark brotli\,1.0.7 result --- top-$1$ agreement of $0$ on three of
four scenarios, despite a moderate pool size --- and the fact that
performance and energy give qualitatively similar but quantitatively
distinct calibration patterns. Aggregating \StudyEnergy with the other tiers
would mix four incompatible sampling mechanisms; we keep them
separate.

\paragraph*{Cost dimensions beyond time.}
Wall-clock measurement time is the canonical operational constraint
and the one we instantiate for \StudyReplication. The x264 $2\times 2$
matrix exercises one non-time cost dimension (output size as cost),
the Linux study exercises another (uncompressed binary size as cost),
and the \StudyEnergy study exercises a third (energy as an additional cost
property alongside performance, with both used as cost and as
target). The qualitative pattern of cost heterogeneity (x264 size:
$\text{CV}=0.04$--$0.31$, $p_{90}/p_{10}=1.12$--$2.21$; Linux:
$\text{CV}\!=\!1.0$--$1.4$, $p_{90}/p_{10}\!=\!4.0$--$6.0$ across
the three releases), calibration miscalibration ($\sim\!\!0.96$--
$1.03$ median on x264 size, strongly strategy-dependent on Linux),
and Pareto-overlap change still hold under the size-based cost
dimensions. The mathematical content of $\TrainCost$ is unchanged
under such substitutions; only the operational interpretation
changes. Replicating the empirical findings under energy-bound or
bandwidth-bound $\mathit{cost}(c)$ is a clean follow-up.

\paragraph*{Implications for prior conclusions.}
The data partly validates and partly contradicts the size-as-cost
abstraction. On systems with low cost heterogeneity (LLVM, Polly), the
abstraction is essentially fine and previously published rankings under
$(n, \MRE)$ transfer to $(\TrainCost, \MRE)$ without rework. On systems
with high heterogeneity (lrzip, BerkeleyDBC) and on x264 scenarios with
target--cost asymmetry, the previously published rankings deserve
re-statement in terms of empirical cost. We recommend that future
configurable-systems sampling studies report both $\SizeCost$ and
$\TrainCost$ alongside any accuracy metric, so that the cost-model
choice is visible rather than hidden. For Pareto-style claims, Li \etal's
methodological guidance is also directly relevant: papers should report the
front comparison that answers their actual decision question, rather than
defaulting to a generic Pareto indicator whose assumptions are left implicit
\cite{li2020pareto}. The same caution extends beyond the classical sampling
papers we revisit. Configurable-systems research often compares methods under
numbers of measurements, evaluations, or count-based budgets because those
quantities are convenient, reproducible, and often directionally correlated
with real cost. Our results show that this convention is not arbitrary, but it
is also not automatically calibrated. The safe interpretation is therefore
simple: count-based findings remain informative about ordering effort, while
any claim about actual resource savings should be read as provisional unless
measured acquisition cost is shown to track the count budget. Papers that
already calibrate counts from measured wall-clock budgets take a stronger first
step, and our results support that direction rather than contradict it
\cite{chen2021multiobj,chen2023aspirations}.

\paragraph*{Beyond configurable software systems.}
Real-runtime ML and HPO benchmarks (\eg LCBench, NAS-Bench-201) are an
interesting cross-domain testbed for the same calibration question. We
deprioritized them in this paper to keep the empirical claim anchored
in configurable software, but they remain a clean candidate for future
cross-domain replication.
